\begin{lemma}[Prescribed Kempe path height bound]
Let $(T,e_0,\varphi)$ be a globally maximum $k$-triple, with first two
vertices $y_0,y_1$.  If
$\alpha\in\overline\varphi(y_0)$ and
$\beta\in\overline\varphi(y_1)$, then the full
$(\alpha,\beta)$-component through $y_0$ has a simple incidence path $P$
from $y_0$ to $y_1$.  Every vertex-node of $P$ belongs to $V(T)$, its
incidence length is divisible by four, and
\[
 \frac{\operatorname{length}(P)}2+1\leq h(G).
\]
Here the quotient converts incidence length to the number of multigraph
edge identities on the path; the additional one converts edge count to
multigraph vertex count.
\end{lemma}
\begin{proof}
Write $X=V(T)$ as in \noderef{TreeSequenceVertexSet}.  The globally maximum hypothesis in
\noderef{GlobalMaximumTriple} says that $T$ is a Tashkinov tree for the
$k$-triple $(G,e_0,\varphi)$ and that $|X|=t(G)$.  It also makes $T$
maximal for this fixed pair in the sense of
\noderef{FixedPairMaximalTashkinovTree}.  Indeed, if a Tashkinov tree $S$
for the same pair had $X\subseteq V(S)$, then the definition
\noderef{TashkinovOrder}, applied using the $k=\chi'(G)-1$ equality in
\noderef{KTriple}, gives $|V(S)|\leq t(G)=|X|$.  Finiteness and the
inclusion therefore give $V(S)=X$.

Apply \noderef{ScheideMaximalKempePath} with $S=T$, indices $0<1$, and
the prescribed colors $\alpha,\beta$.  It gives
$\alpha\ne\beta$ and a simple path $P$ in the two-color incidence graph
\noderef{TwoColorIncidenceGraph}, from the vertex-node $y_0$ to the
vertex-node $y_1$, whose support is exactly the full component
\noderef{TwoColorComponent} through $y_0$.  That result also gives
$4\mid\operatorname{length}(P)$ and places every vertex-node of $P$ in
$X$.  Notice that the path alternates between vertex-nodes and individual
edge-nodes.  Thus an incidence length of $2m$ represents exactly $m$
multigraph edge identities, even when two such identities have the same
unordered pair of endpoints.

It remains to compare this prescribed path with $h(G)$.  Reverse $P$ and
write its multigraph vertices and edge identities as
\[
 y_1,f_m,z_{m-1},f_{m-1},\ldots,z_1,f_1,y_0,
 \qquad m=\operatorname{length}(P)/2.
\]
The edge identities $f_1,\ldots,f_m$ are pairwise distinct because $P$
is a simple incidence path.  None equals $e_0$: by
\noderef{TwoColorIncidenceGraph} every displayed edge-node is colored
$\alpha$ or $\beta$, whereas the $k$-triple clause
\noderef{KTriple} says that $e_0$ is uncolored.  The endpoints $y_0,y_1$
are distinct, so the positive path length together with
$4\mid\operatorname{length}(P)$ gives $m\geq2$.  Form an initial tree
$U$ by starting with
$y_0,e_0,y_1$, then adjoining successively
$f_m,z_{m-1},\ldots,f_2,z_1$; reserve $f_1$ as the closing edge from
$z_1$ to $y_0$.  This construction uses the actual edge-nodes of $P$, so
parallel edges are never identified.  Simplicity of $P$, injectivity in
\noderef{TreeSequence}, and looplessness give distinct displayed vertices,
distinct displayed tree-edge identities, and the required attachment
equalities.

The path colors alternate between $\alpha$ and $\beta$.  Every colored
tree edge added after $e_0$ therefore has a missing-color witness among
$y_0,y_1$, by the two prescribed hypotheses and
\noderef{MissingColors}.  Hence $U$ satisfies
\noderef{TashkinovTree}.  The same fixed-pair maximality established in
the first paragraph now lets
\noderef{ScheideMaximalExtendsEveryTree} extend $U$, without changing its
ordered initial segment, to a Tashkinov tree $T'$ with $V(T')=X$.  The
initial-segment assertion is entrywise equality in the sense of
\noderef{TreeSequenceInitialSegment}.  The $k$-triple is unchanged, and
$|V(T')|=|X|=t(G)$, so $T'$ is again globally
maximum.

The retained initial segment, the reversed incidence path, and the
closing identity $f_1$ define normal data in the sense of
\noderef{NormalGlobalMaximumTriple}.  Its chosen height is the number of
multigraph vertices on $P$, namely
$m+1=\operatorname{length}(P)/2+1$.  The support clauses hold because the
incidence path is simple: its vertex-nodes are precisely the $m+1$
vertices of the initial segment, and its edge-nodes are precisely the
$m-1$ added tree identities together with the closing identity.

Finally this normal globally maximum triple is one of the witnesses over
which \noderef{ScheideGraphHeight} takes its finite maximum.  Its height
is at most $|V(T')|\leq |V(G)|$, so it lies in the defining finite range.
Consequently
\[
 \operatorname{length}(P)/2+1=m+1\leq h(G),
\]
which proves the asserted bound for the originally prescribed pair.
\end{proof}
